\section{Grammar Generation Ability of LLMs}
In this section, we describe the dedicated dataset constructed to evaluate the ability of LLMs in grammar generation, introduce 6 metrics we use in evaluation, explain experiments, and analyze results in detail. We detail each in the following subsections.

\subsection{Dataset}
To evaluate the capacity of LLMs for grammar generation, we present a dedicated dataset. 

During dataset construction, for each $k \in \{1,2,\dots,9\}$, we prompted GPT-4o~\cite{openai2024gpt4technicalreport} to generate 10 distinct reference grammars, where each reference grammar $\bnfref$ has $k$ nonterminal symbols and thus $|R| = k$. This gives a total of $90$ reference grammars.
% consists of exactly $k$ grammar
% edge elements. 
We used $\bnfref$ to prompt GPT-4o to produce 6 different challenges with each challenge consisting of a set of positives $\mathcal{P}$ and negatives $\mathcal{N}$ where $|\mathcal{P}|=3$ and $|\mathcal{N}|=3$, in a way that $\mathcal{P} \subseteq \languageOf{\bnfref}$ and $\mathcal{N} \cap \languageOf{\bnfref} = \emptyset$. However, GPT-4o often failed to produce challenges with valid reference grammars, positives, and negatives as $k$ increased. We manually corrected erroneous reference grammars, positives, and negatives by using a BNF parser which takes a grammar and outputs whether a grammar is in valid BNF, and whether positives are accepted, and negatives are rejected\footnote{Refer to Appendix~\ref{app:dataset-construction} for the details of dataset construction.}.

Following this process, we obtained a dataset of 540 challenges, each consisting of 3 positives and 3 negatives.
Figure~\ref{fig:bnf_generation} shows an example challenge and a corresponding solution. 

% consider a challenge with $\mathcal{P}$ as:
% \begin{center}
% \begin{minipage}{0.35\linewidth}
% \begin{verbatim}
% add(1,2,3)
% merge(x,y)
% fibonacci(9)
% \end{verbatim}
% \end{minipage}
% \end{center}
% and $\mathcal{N}$ as:
% \begin{center}
% \begin{minipage}{0.35\linewidth}
% \begin{verbatim}
% add(1 2 3)
% merge[x,y]
% fibonacci 9
% \end{verbatim}
% \end{minipage}
% \end{center}
% An example valid $\bnf$ such that $\mathcal{P} \subseteq \languageOf{\bnf}$ and $\mathcal{N} \cap \languageOf{\bnf} = \emptyset$ is as follows:
% \begin{center}
% \begin{minipage}{0.92\linewidth}
% \begin{verbatim}
% <stmt> ::= <func> "(" <args> ")"
% <args> ::= <expr> | <expr> "," <args>
% <expr> ::= <char> | <number>
% <func> ::= <char> <func> | <char>
% <char> ::= "a" | ... | "z"
% <number> ::= "0" | ... | "9"
% \end{verbatim}
% \end{minipage}
% \end{center}

\subsection{Metrics}
Let $C$ be a set of $N$ challenges where each is a tuple 
$(\bnfref,\mathcal{P},\mathcal{N}, \guess)$ consisting of a reference grammar, and a set of positive examples and negative examples, respectively, and the corresponding candidate grammar $\guess$ generated by an LLM. We evaluate the quality of generated grammars using 6 key metrics:

\paragraph{Syntax Correctness ($\sx$)}
The syntax correctness metric $\sx$ quantifies the proportion of guesses that conform to the valid BNF syntax defined in Section~\ref{sec:bnf_background}.
% A grammar generated by the LLM, $\guess$, is deemed syntactically correct if and only if it is successfully parsed by a BNF parser. 
% If the grammar is syntactically correct,we construct the corresponding grammar edge set $\mathcal{E}^\mathcal{G}$; otherwise, the grammar is rejected, and we have $\mathcal{E}^\mathcal{G}=\emptyset$
We define an indicator function as:
\begin{align*}
\mathbb{I}_{SX}(\guess) &=
\begin{cases}
1 & \text{if $valid(\guess)$,}\\
0 & \text{otherwise.}
\end{cases}
\end{align*}

\noindent $SX(C)$ is defined as:
$
% \mathrm{SX}(C) &= 
\frac{1}{N}\sum_{i=1}^N \mathbb{I}_{SX}(\guess_{i}).
$

\paragraph{Semantic Correctness ($\se$)} $\se$ captures the proportion of guesses that are semantically correct.
%
We define an indicator function as follows, noting that if $\guess$ is not in valid BNF, then $\languageOf{\guess} = \emptyset$:
\begin{align*}
\mathbb{I}_{SE}(\guess,\mathcal{P},\mathcal{N}) &=
\begin{cases}
1 & \text{if $\mathcal{P} \subseteq \languageOf{\guess} \,\wedge$} \\
& \mathcal{N} \cap \languageOf{\guess} = \emptyset \\
0 & \text{otherwise.}
\end{cases}
\end{align*}


% \begin{align*}
% \mathbb{I}_{SE}(\guess,\mathcal{P},\mathcal{N}) &=
% \begin{cases}
% 1 & \text{if } \big(\text{($\mathcal{E}^{\mathcal{G}}\neq \phi$)}\\
% & \quad \wedge\; (\mathcal{P}_i \subseteq \mathcal{L}^\mathcal{G}) \\
% & \quad \wedge\; (\mathcal{N}_i \centernot{\cap} \mathcal{L}^\mathcal{G})\big)\text{,} \\
% 0 & \text{otherwise.}
% \end{cases}
% \end{align*}
% meaning if only if a given grammar has correct syntax and it can accept all positive examples and reject all negative examples, it has correct semantics.

\noindent$SE(C)$ is given by:
$
\frac{1}{N}\sum_{i=1}^N \mathbb{I}_{SE}(\guess_{i},\mathcal{P}_i,\mathcal{N}_i)
$.

\paragraph{Estimating Grammar Quality}
Given a set of positive and negative examples, there are many possible valid and semantically correct solutions that are undesirable. For instance, the following grammar would be an undesirable solution for Figure~\ref{fig:bnf_generation}. Here the $\languageOf{\guess} = |P|$, and is overfitted to the examples:

\vspace{1em}
\begin{minipage}{0.92\linewidth}
\begin{verbatim}

<stmt> ::= "add(1,2,3)" | "merge(x,y)" |
           "fibonacci(9)"
\end{verbatim}
\end{minipage}
\vspace{1em}


Equally undesirable is a grammar that over-generalizes from the examples, and defines a significantly larger language than the reference grammar, because this is highly likely to contain invalid strings. 
%
There are no standard metrics to measure over-fitting or over-generalization in grammar generation, so we devise $4$ metrics based around the number of production rules used in parsing the positive examples. 

First, let $\Pi_\mathcal{P} \subseteq \Pi$ be the set of production rules that are used in the left-most derivations of all positive examples in $\mathcal{P}$. That is, the set of rules in $\Pi$ which occur in a sequence of rules $S \rightarrow \alpha_1 \rightarrow \ldots \rightarrow \alpha_n \rightarrow  p $ where $p \in \mathcal{P}$, and all rules expand the left-most non-terminal in $\alpha_1, \ldots, \alpha_n$. 

We report metrics across only \emph{solved} challenges, i.e., a challenge where $\guess$ is syntactically and semantically correct. Our four metrics are defined as follows\footnote{Refer to Appendix~\ref{appendix:grammar_quality_metrics} for the details of formal definitions.}:
\begin{squishitem}
    \item $\diff(C)$, calculates the difference between the number of production rules in $\guess$ used in parsing the positive examples and the number of production rules used by $\bnfref$ for a given challenge, i.e., $|\Pi^{ref}_\mathcal{P}| - |\Pi^*_\mathcal{P}|$. A grammar that uses substantially fewer production rules has probably overfitted to the examples, and a grammar that uses substantially more production rules may have over-generalized. We report the average of $\diff$ across all solved challenges as $\diff^\diamond$.
    \item $\of$ estimates over-fitting by counting the percentage of solved challenges on which $\guess$ uses fewer than half the number of production rules used by $\bnfref$, i.e., the number of times $|\Pi^{ref}_\mathcal{P}| - |\Pi^*_\mathcal{P}| > \frac{|\Pi^{ref}_\mathcal{P}|}{2}$. 
    \item $\og$ estimates over-generalization by counting the percentage of solved challenges on which $\guess$ uses more than half the number of production rules used by $\bnfref$, i.e.,  $|\Pi^{ref}_\mathcal{P}| - |\Pi^*_\mathcal{P}| < -\frac{|\Pi^{ref}_\mathcal{P}|}{2}$.
    \item $\tu$ calculates the proportion of production rules that are used in parsing the positive examples for a given challenge, i.e., $\frac{|\Pi^*_\mathcal{P}|}{\Pi^*}$, indicating the utility. We report the average of $\tu$ across all solved challenges as $TU^\diamond$.
\end{squishitem} 

% These metrics are used only for syntactically and semantically correct grammars, and are thus calculated over $C' \subseteq C$, where $C'$ consists of  $k$ challenges from $C$ for which the LLM has generated a syntactically and semantically correct grammar.

% $FGI$ measures the fitting and generalization of grammars, consisting of three metrics $FGI^{Diff}$, $FGI^{OF}$, and $FGI^{OG}$. They are used only for grammars that $\mathcal{E}^{\mathcal{G}^{llm}_i}\neq \emptyset$.


% First define $\Pi_\mathcal{P} \subseteq \Pi$ to be the set of production rules that are used in the left-most derivations of all positive examples in $\mathcal{P}$. That is, the set of rules in $\Pi$ which occur in a sequence of rules $S \rightarrow \alpha_1 \rightarrow \ldots \rightarrow \alpha_n \rightarrow \rightarrow p $ where $p \in \mathcal{P}$, and all rules expand the left-most non-terminal in $\alpha_1, \ldots, \alpha_n$. 


% % First define a $\mathrm{ProdUsed}$ function for measuring the number of production rules used during derivation for a grammar $\mathcal{G}$: 
% % \[
% % \mathrm{ProdUsed}(\mathcal{E}^{\mathcal{G}}, \mathcal{P})
% % \;=\;
% % \Bigl|\{\,a \in R
% % \;\bigm|\;
% % S \overset{*}{\underset{a}{\Rightarrow}} p
% % \}\Bigr|.
% % \]
% % where $p \in \mathcal{P}$ , and $S$ derives $p$ in zero or more steps using production rule $a \in R$. 

% Let $\Pi^*$ be the production rules in $\guess$ and $\Pi^{ref}$ be the production rules in $\bnfref$. Let us define $\textit{Diff}(\bnfref,\guess) =  |\Pi^{ref}_\mathcal{P}| - |\Pi^*_\mathcal{P}|.$
% %
% %
% The average difference in production rules used over the whole set of solved challenges, $C'$, is given by: 
% $$
% \diff(C') = \frac{1}{k} \sum_{i=1}^k \mathit{Diff}(\bnfref_i,\guess_i).
% $$


% % We define the difference of the number of used production rules between two grammars as:
% % \[
% % \mathrm{PUDiff}(\mathcal{E}^{\mathcal{G}_{(1)}},\mathcal{E}^{{\mathcal{G}}_{(2)}},\mathcal{P}) = \mathrm{PU_{(1)}} - \mathrm{PU_{(2)}}
% % \]
% % where $\mathrm{PU_{(1)}} = \mathrm{ProdUsed}(\mathcal{E}^{\mathcal{G}_{(1)}},\mathcal{P})$ and $\mathrm{PU_{\mathcal{G}^2}} = \mathrm{ProdUsed}(\mathcal{E}^{\mathcal{G}_{(2)}},\mathcal{P})$ to indicate the difference of the number of used production rules between $\mathcal{G}_{(1)}$ and $\mathcal{G}_{(2)}$.

% % The $FGI^{Diff}$ is given as:
% % \begin{align*}
% %     \mathrm{FGI^{Diff}}(C) = \frac{1}{N} \sum_{i=1}^N \mathrm{PUDiff}(\mathcal{E}^{\mathcal{G}_{ref}}_i,\mathcal{E}^{\mathcal{G}_{llm}}_i,\mathcal{P}_i).
% % \end{align*}
% % to measure the average difference of used production rules between reference and generated grammars.

% % Let $\mathrm{PUDiff}^{(1)}_{(2)}$ be $\mathrm{PUDiff}(\mathcal{E}^{\mathcal{G}_{(1)}},\mathcal{E}^{\mathcal{G}_{(2)}},\mathcal{P})$ and $\mathrm{ProdUsed^{(1)}}$ be $\mathrm{ProdUsed(\mathcal{E}^{\mathcal{G}_{(1)}},\mathcal{P})}$. 

% \noindent We define two indicator functions, which indicates when a grammar uses substantially fewer rules than the reference grammar, and substantially more rules than the reference grammar:
% % \begin{align*}
% % \mathbb{I}_{OF}(\mathcal{E}^{\mathcal{G}_{(1)}},\mathcal{E}^{\mathcal{G}_{(2)}},\mathcal{P}) &= 
% % \begin{cases}
% % 1 & \text{if } (\mathrm{PUDiff}^{(1)}_{(2)} \\ & \quad > \frac{\mathrm{ProdUsed^{(1)}}}{2}) \\
% % 0 & \text{otherwise.}
% % \end{cases}
% % \end{align*}
% \begin{align*}
% \mathbb{I}_{OF}(\bnfref, \guess,\mathcal{P}) &= 
% \begin{cases}
% 1 & \text{if } |\Pi^{ref}_\mathcal{P}| - |\Pi^*_\mathcal{P}|
% % \\ & \quad
% > \frac{|\Pi^{ref}_\mathcal{P}|}{2} \\
% 0 & \text{otherwise.}
% \end{cases}
% \end{align*}
% \begin{align*}
% \mathbb{I}_{OG}(\bnfref, \guess,\mathcal{P}) &= 
% \begin{cases}
% 1 & \text{if } |\Pi^{ref}_\mathcal{P}| - |\Pi^*_\mathcal{P}|
% % \\ & \quad
% < -\frac{|\Pi^{ref}_\mathcal{P}|}{2} \\
% 0 & \text{otherwise.}
% \end{cases}
% \end{align*}



% The metric to estimate whether the generated grammars overfit the examples, is then given by
% $$
%  \of(C') = \frac{1}{k}\sum_{i=1}^{k} \mathbb{I}_{OF}(\bnfref_i,\guess_i, \mathcal{P}_i).
% $$


% % We also define an indicator function:
% % \begin{align*}
% % \mathbb{I}_{OG}(\mathcal{E}^{\mathcal{G}_{(1)}},\mathcal{E}^{\mathcal{G}_{(2)}},\mathcal{P}) &= 
% % \begin{cases}
% % 1 & \text{if } (\mathrm{PUDiff}^{(1)}_{(2)} \\ &\quad < -\frac{\mathrm{ProdUsed^{(1)}}}{2}) \\
% % 0 & \text{otherwise.}
% % \end{cases}
% % \end{align*}
% The metric to estimate whether the generate grammars over generalize the examples, $FGI^{OG}$, is given as:
% \begin{align*}
%     \og(C') = \frac{1}{k}\sum_{i=1}^{k}  \mathbb{I}_{OG}(\bnfref_i, \guess_i, \mathcal{P}_i).
% \end{align*}

% \subsection{TU Index ($\tu$)} 
% Given a set of positives and a guess $G^{*}$, the percentage of $|\Pi^*|$ taken up by $|\Pi^*_\mathcal{P}|$ indicating the utility of $G^{*}$. We define a metric $\tu$ to measure it:
% $TU = \frac{|\Pi^*_\mathcal{P}|}{|\Pi^*|}$ for which lower $\tu$ indicates a bunch of irrelevant or nonsensical production rules of $G^{*}$ while higher $\tu$ indicates the opposite. Same as $FG$, $\tu$ are used only for syntactically and semantically correct grammars.

\subsection{Baselines}
To establish baselines, we adopted two approaches, Direct Prompting (DP) and Optimization of the BNF Parser for LLM-Friendly Feedback (OPF).

In DP, we directly prompted LLMs with positive and negative examples, asking them to produce grammars that accept all positives and reject all negatives\footnote{Refer to Appendix~\ref{appendix:DP} for the details of DP.}.

In OPF, inspired by Reflexion~\cite{shinn2023reflexion} and Self-Refine~\cite{madaan2023selfrefine}, we further optimized the BNF parser by enabling it to provide more LLM-friendly error messages as feedback, aiming to improve the performance of grammar generation for LLMs\footnote{Refer to Appendix~\ref{appendix:OPF} for the details of OPF.}. 

\subsection{Experiment Settings}
For a comprehensive evaluation, we selected a total of 8 LLMs, ensuring a diverse selection of both open- and closed-source LLMs, along with LLMs with varying parameter sizes and LLMs specifically designed for code generation. 

Specifically, we selected two closed-source LLMs, \textit{GPT-4o}~\cite{openai2024gpt4technicalreport} and \textit{GPT-3.5-Turbo}~\cite{brown2020languagemodelsfewshotlearners}. For the other 6 open-source LLMs, we note them in the notation: \text{\{\textit{model\_name}\}\textit{:}\{\textit{parameter\_size}\}\textit{-}\{\textit{model\_type}\}.} We selected the following LLMs: \textit{Llama3:70b-Instruct}~\cite{grattafiori2024llama3herdmodels}, \textit{Qwen:72b-Chat}~\cite{bai2023qwentechnicalreport}, \textit{Gemma2:27b-Instruct}~\cite{gemmateam2024gemma2improvingopen}, \textit{Mistral:7b-Instruct}~\cite{jiang2023mistral7b}, \textit{Codestral:22b}~\cite{mistral2024codestral}, and \textit{Starcoder2:15b-Instruct}~\cite{lozhkov2024starcoder2stackv2}. Among the selected LLMs, \textit{Starcoder2:15b-Instruct} and \textit{Codestral:22b} are LLMs for code generation.

For the DP baseline, we set the temperature to 0 and the maximum token to 2000.  

For the OPF baseline\footnote{Refer to Appendix~\ref{appendix:OPF} for the reason of setting the temperature greater than 0 for OPF.}, we set the temperature to 0.3, the maximum token to 2000, and the maximum number of turns, \textit{max\_turns}, to 5.


\subsection{Results \& Analysis}
The results of $\sx$ and $\se$ for the 8 LLMs are presented in Table~\ref{tab:results_for_all}. We observed that \textit{GPT-4o}, \textit{GPT-3.5-Turbo}, and \textit{Gemma2:27b-Instruct} achieve relatively high $\sx$, while the other LLMs generally fare worse in DP baseline. Compared to the DP baseline, applying OPF leads to a significant enhancement in $\sx$ for both \textit{Mistral:7b-Instruct} and \textit{Codestral:22b}, yielding an 18\% improvement for each of them, suggesting that parser feedback can help increase $\sx$. Nevertheless, for most LLMs, OPF yields only slight gains in $\sx$. Furthermore, it is worth noting that for \textit{Starcoder2:15b-Instruct}, $\sx$ decreases by 16\% after applying OPF. This indicates that feedback from the parser can sometimes lead to performance degradation if LLMs fail to interpret the feedback correctly. 

Although several LLMs attain high $\sx$, their $\se$ remains low in DP, with an average of 39\%. For example, \textit{GPT-3.5-Turbo} has 94\% $\sx$ but only 37\% $\se$. With OPF, LLMs like \textit{Mistral:7b-Instruct} and \textit{Codestral:22b} achieve notable $\se$ improvement with enhanced $\sx$, illustrating that improved $\sx$ can positively influence $\se$. However, for most LLMs, OPF yields only marginal $\se$ gains, particularly in the case that their $\sx$ is already high and OPF fails to contribute significant $\sx$ improvement.

Noticeably, although OPF helps improve $\sx$, a significant gap still persists between $\sx$ and $\se$ for most LLMs. For example, even after OPF, \textit{GPT-3.5-Turbo} maintains a 57\% gap between $\sx$ and $\se$, and \textit{Starcoder2:15b-Instruct} has a 40\% gap. Since the ultimate objective is to improve $\se$, a more advanced approach is needed to address this limitation.

Furthermore, as other four metrics shown in Table~\ref{tab:other_metrics}, on average, for most of LLMs, the lower $\diff^\diamond$, $\of$, and $\og$ indicate negligible over-fitting or over-generalization issues, and the higher $\tu^\diamond$ means LLMs are not predisposed to generate irrelevant production rules. Nevertheless, as the number of production rules increases, $\diff^\diamond$ and $\of$ exhibit slight increases, indicating a tendency toward mild over-fitting, while no significant over-generalization issues are observed~\footnote{Refer to Appendix~\ref{appendix:fgi_results} and~\ref{appendix:tu_results} for more details.}.

In addition to the quantitative metrics, we also conducted a qualitative analysis by examining the generated grammars. We compared the grammars containing syntactic errors generated in the DP but corrected after the OPF. We primarily found four significant issues causing lower $SX$: unsupported symbols such as injecting quantifiers like ``*'' and ``?'' or character classes like ``[a-z]'', erroneously introduced and misplaced brackets such as wrapping two terminals with round and square brackets, failing to wrap non-terminal with angle brackets, and omission the separators ``|'' between sentential forms. While the the first two issues are rampant across most LLMs, the third issue is mainly found in \textit{Mistral:7b-Instruct} and the last one is sporadic. For most LLMs, OPF can occasionally mitigate these issues, but insignificantly. Furthermore, we also observed LLMs show an ability to recognize keywords in the examples, treating them as complete terminals rather than decomposing them into multiple terminals as individual characters. For example, they treat ``if'' and ``SELECT'' as complete terminals rather than splitting them into multiple characters. This may be benefited from the common sense acquired by LLMs from the corpus~\cite{llm_few_shot}. 

 % we have also demonstrated the results of fitting and generalization metrics $FGI^{Diff}$, $FGI^{OF}$, $FGI^{OG}$ in Table~\ref{tab:fgi_diff},~\ref{tab:fgi_of}, and ~\ref{tab:fgi_og} in Appendix~\ref{appendix:fgi_results}, respectively.

% Syntax and semantic correctness results of all LLMS with different methods
% \begin{table*}
% \centering
% \begin{tabularx}{\textwidth}{l X X X X X l}
% \toprule
% \textbf{Models} & \textbf{$SX_{DP}$} & \textbf{$SX_{OPF}$} & \textbf{$SX_{\NAME}$} & \textbf{$SE_{DP}$} & \textbf{$SE_{OPF}$} & \textbf{$SE_{\NAME}$}\\
% \midrule
% GPT-4o
% & 93 
% & \textbf{97}~\textcolor{blue}{\scriptsize$\uparrow$4} 
% & 96~\textcolor{blue}{\scriptsize$\uparrow$3}~\textcolor{red}{\scriptsize$~\downarrow$1} 
% & 84 
% & 85~\textcolor{blue}{\scriptsize$\uparrow$1}
% & \textbf{93}~\textcolor{blue}{\scriptsize$\uparrow$9}~\textcolor{red}{\scriptsize$~\uparrow$8}  \\
% GPT-3.5-Turbo
% & 94 
% & 95~\textcolor{blue}{\scriptsize$\uparrow$1} 
% & \textbf{99}~\textcolor{blue}{\scriptsize$\uparrow$5}~\textcolor{red}{\scriptsize$~\uparrow$4} 
% & 37 
% & 38~\textcolor{blue}{\scriptsize$\uparrow$1}
% & \textbf{61}~\textcolor{blue}{\scriptsize$\uparrow$24}~\textcolor{red}{\scriptsize$\uparrow$23}  \\
% Llama3:70b-Instruct
% & 57 
% & 61~\textcolor{blue}{\scriptsize$\uparrow$4}
% & \textbf{75}~\textcolor{blue}{\scriptsize$\uparrow$18}~\textcolor{red}{\scriptsize$\uparrow$14} 
% & 41 
% & 42~\textcolor{blue}{\scriptsize$\uparrow$1} 
% & \textbf{61}~\textcolor{blue}{\scriptsize$\uparrow$20}~\textcolor{red}{\scriptsize$\uparrow$19}  \\
% Qwen:72b-Chat
% & 47 
% & 49~\textcolor{blue}{\scriptsize$\uparrow$2}
% & \textbf{76}~\textcolor{blue}{\scriptsize$\uparrow$29}~\textcolor{red}{\scriptsize$\uparrow$27}
% & 20 
% & 21~\textcolor{blue}{\scriptsize$\uparrow$1}
% & \textbf{38}~\textcolor{blue}{\scriptsize$\uparrow$18}~\textcolor{red}{\scriptsize$\uparrow$17} \\
% \textit{Mistral:7b-Instruct }
% & \phantom{0}1 
% & \textbf{19}~\textcolor{blue}{\scriptsize$\uparrow$18} 
% & \phantom{0}1~\textcolor{blue}{~\scriptsize$-$} ~\textcolor{red}{\scriptsize$\downarrow$18}
% & \phantom{0}0 
% & \phantom{0}\textbf{8}~\textcolor{blue}{\scriptsize$\uparrow$8}
% & \phantom{0}1~\textcolor{blue}{\scriptsize$\uparrow$1}~\textcolor{red}{\scriptsize$~\downarrow$7}\\
% Gemma2:27b-Instruct
% & 91 
% & 92~\textcolor{blue}{\scriptsize$\uparrow$1} 
% & \textbf{98}~\textcolor{blue}{\scriptsize$\uparrow$7}~\textcolor{red}{\scriptsize$~\uparrow$6}
% & 56 
% & 57~\textcolor{blue}{\scriptsize$\uparrow$1} 
% & \textbf{79}~\textcolor{blue}{\scriptsize$\uparrow$23}~\textcolor{red}{\scriptsize$\uparrow$22} \\
% Starcoder2:15b-Instruct
% & 76 
% & 60~\textcolor{blue}{\scriptsize$\downarrow$16} 
% & \textbf{98}~\textcolor{blue}{\scriptsize$\uparrow$22}~\textcolor{red}{\scriptsize$\uparrow$38}
% & 30 
% & 20~\textcolor{blue}{\scriptsize$\downarrow$10} 
% & \textbf{44}~\textcolor{blue}{\scriptsize$\uparrow$14}~\textcolor{red}{\scriptsize$\uparrow$24} \\
% Codestral:22b
% & 53 
% & 71~\textcolor{blue}{\scriptsize$\uparrow$18} 
% & \textbf{80}~\textcolor{blue}{\scriptsize$\uparrow$27}~\textcolor{red}{\scriptsize$\uparrow$9}
% & 44 
% & 52~\textcolor{blue}{\scriptsize$\uparrow$8} 
% & \textbf{67}~\textcolor{blue}{\scriptsize$\uparrow$23}~\textcolor{red}{\scriptsize$\uparrow$15} \\
% \bottomrule
% \end{tabularx}
% \caption{Results for LLM grammar generation are presented as percentages for syntax and semantic correctness (\%). For each LLM, the best syntax and semantic score are highlighted with bold font. Blue arrows \textcolor{blue}{$\uparrow\downarrow$} represent performance differences relative to the DP baseline, while red arrows \textcolor{red}{$\uparrow\downarrow$} indicate differences relative to the OPF baseline.}
% \label{tab:results_for_all}
% \end{table*}

\begin{table*}
\centering
\begin{tabularx}{\textwidth}{l *{6}{>{\centering\arraybackslash}X}}
\toprule
{\textbf{Models}} & \multicolumn{3}{c}{\textbf{Syntax Correctness ($\sx$)}} & \multicolumn{3}{c}{\textbf{Semantics Correctness($\se$)}} \\
\cmidrule(lr){2-4} \cmidrule(lr){5-7}
 & DP & OPF & \textbf{\NAME} & DP & OPF & \textbf{\NAME} \\
\midrule
GPT-4o 
& 93 
& \textbf{97}~\textcolor{blue}{\scriptsize$\uparrow$4} 
& 96~\textcolor{blue}{\scriptsize$\uparrow$3}~\textcolor{red}{\scriptsize$~\downarrow$1} 
& 84 
& 85~\textcolor{blue}{\scriptsize$\uparrow$1}
& \textbf{93}~\textcolor{blue}{\scriptsize$\uparrow$9}~\textcolor{red}{\scriptsize$~\uparrow$8}  \\
GPT-3.5-Turbo 
& 94 
& 95~\textcolor{blue}{\scriptsize$\uparrow$1} 
& \textbf{99}~\textcolor{blue}{\scriptsize$\uparrow$5}~\textcolor{red}{\scriptsize$~\uparrow$4} 
& 37 
& 38~\textcolor{blue}{\scriptsize$\uparrow$1}
& \textbf{61}~\textcolor{blue}{\scriptsize$\uparrow$24}~\textcolor{red}{\scriptsize$\uparrow$23}  \\
Llama3:70b-Instruct 
& 57 
& 61~\textcolor{blue}{\scriptsize$\uparrow$4}
& \textbf{75}~\textcolor{blue}{\scriptsize$\uparrow$18}~\textcolor{red}{\scriptsize$\uparrow$14} 
& 41 
& 42~\textcolor{blue}{\scriptsize$\uparrow$1} 
& \textbf{61}~\textcolor{blue}{\scriptsize$\uparrow$20}~\textcolor{red}{\scriptsize$\uparrow$19}  \\
Qwen:72b-Chat 
& 47 
& 49~\textcolor{blue}{\scriptsize$\uparrow$2}
& \textbf{76}~\textcolor{blue}{\scriptsize$\uparrow$29}~\textcolor{red}{\scriptsize$\uparrow$27}
& 20 
& 21~\textcolor{blue}{\scriptsize$\uparrow$1}
& \textbf{38}~\textcolor{blue}{\scriptsize$\uparrow$18}~\textcolor{red}{\scriptsize$\uparrow$17} \\
Mistral:7b-Instruct
& \phantom{0}1 
& \textbf{19}~\textcolor{blue}{\scriptsize$\uparrow$18} 
& \phantom{0}1~\textcolor{blue}{\scriptsize$-$}~\textcolor{red}{\scriptsize$\downarrow$18}
& \phantom{0}0 
& \phantom{0}\textbf{8}~\textcolor{blue}{\scriptsize$\uparrow$8}
& \phantom{0}1~\textcolor{blue}{\scriptsize$\uparrow$1}~\textcolor{red}{\scriptsize$\downarrow$7}\\
Gemma2:27b-Instruct 
& 91 
& 92~\textcolor{blue}{\scriptsize$\uparrow$1} 
& \textbf{98}~\textcolor{blue}{\scriptsize$\uparrow$7}~\textcolor{red}{\scriptsize$\uparrow$6}
& 56 
& 57~\textcolor{blue}{\scriptsize$\uparrow$1} 
& \textbf{79}~\textcolor{blue}{\scriptsize$\uparrow$23}~\textcolor{red}{\scriptsize$\uparrow$22} \\
Starcoder2:15b-Instruct 
& 76 
& 60~\textcolor{blue}{\scriptsize$\downarrow$16} 
& \textbf{98}~\textcolor{blue}{\scriptsize$\uparrow$22}~\textcolor{red}{\scriptsize$\uparrow$38}
& 30 
& 20~\textcolor{blue}{\scriptsize$\downarrow$10} 
& \textbf{44}~\textcolor{blue}{\scriptsize$\uparrow$14}~\textcolor{red}{\scriptsize$\uparrow$24} \\
Codestral:22b 
& 53 
& 71~\textcolor{blue}{\scriptsize$\uparrow$18} 
& \textbf{80}~\textcolor{blue}{\scriptsize$\uparrow$27}~\textcolor{red}{\scriptsize$\uparrow$9}
& 44 
& 52~\textcolor{blue}{\scriptsize$\uparrow$8} 
& \textbf{67}~\textcolor{blue}{\scriptsize$\uparrow$23}~\textcolor{red}{\scriptsize$\uparrow$15} \\
\bottomrule
\end{tabularx}
\caption{The results of syntax and semantic correctness for LLMs grammar generation are presented as percentages (\%). For each LLM, the best syntax and semantic correctness are highlighted with bold font. Blue arrows \textcolor{blue}{$\uparrow\downarrow$} represent performance differences relative to the DP baseline, while red arrows \textcolor{red}{$\uparrow\downarrow$} indicate differences relative to the OPF baseline.}
\label{tab:results_for_all}
\end{table*}